\section{Rozwój}
\label{sec:rozwoj}
W toku dotychczasowych prac nad stanowiskiem 
testowym \cite{ESAPAC2026_HIL,skromak2026hilall} zidentyfikowano 
kilka aspektów do poprawy w działaniu stanowiska, które
nie zostały zrealizowane w ramach tej pracy. 
Te aspekty to:
\begin{itemize}
    \item Możliwość wprowadzenia korekty mierzonych wychyleń sterów.
    \item Zwiększenie częstotliwości wysyłania danych przez \textit{CASKET}.
\end{itemize}

\subsection{Wbudowana w \textit{CASKET} korekta wychyleń sterów}
\label{subsec:rozwoj-korekta}
Obecnie użycie bezkontaktowych czujników magnetycznych sprawia,
że pomiar kąta wykonywany przez \textit{CASKET} jest niedokładny
i zależy od względnego ustawienia czujnika i magnesu. 
zmitygowanie tego problemu polega na iteracyjnych korekcjach ustawienia
czujnika względem magnesu. Podczas projektowania \textit{CASKET}a 
została przewidziana taka ewentualność - jest to jedna z przyczyn 
zastosowania stolików ze śrubami mikrometrycznymi. Proces 
ten, o ile skuteczny i jednorazowy dla jednego włożenie układu sterowania w urządzenie pomiarowe, 
o tyle znacznie wydłuża czas przygotowania do pierwszego testu, 
a jego wynik w dużej mierze zależy od operatora wprowadzającego poprawki.
Pożądane jest więc wprowadzenie automatyzacji procesu. 
Korekcja względnego ustawienia magnesu i czujnika wciąż jest 
wskazana, jednak powinna być wspomagana przez poprawki wprowadzone
w oprogramowaniu stanowiska.
Zostanie stworzony skrypt, którego zadaniem będzie automatyczne
wygenerowanie poprawek pomiarowych dla \textit{CASKET}a.
Jego działanie będzie polegało na wychyleniu sterów na trzy pozycje:
2 skrajne i zerową, sczytanie pomiarów \textit{CASKET}a
i następnym takim skorygowaniu pomiarów \textit{CASKET}a, by 
te punkty leżały możliwie blisko prostej \(y = x\).
Planuje się do tego zastosować przekształcenie liniowe.
Należy pamiętać o wpływie luzów występujących w systemie 
\todo[inline]{Wymyślić czy jest coś lepszego niż średnia z ruchu w jedną i drugą stronę}
Ta automatyczna procedura może też stanowić ulepszenie 
względem obecnego ustawiania \textit{zera} czujników, mitygując
problem luzów w występujący w obecnie używanej procedurze.

\subsection{Zwiększenie częstotliwości \textit{CASKET}a}
\label{subsec:rozwoj-zwiekszzenie-czestotliwosci}
O ile możliwe jest prowadzenie pomiarów z częstotliwością \SI{1600}{\hertz},
to w toku testów zdiagnozowano czasową utratę sygnału z \textit{CASKET}a.
Problem rozwiązano zmniejszając częstotliwość wysyłanego wychylenia 
do \SI{800}{\hertz}. Jest to wciąż duża częstotliwość i zmniejszenie
jej nie wpływa znacząco na symulację.
\todo[inline]{Policzyć ile czasu przy obecnej częstotliwośc}
Zwiększenie częstotliwości wysyłanych danych możliwe będzie 
gdy użyte zostanie \textit{Direct Memory Acces} (DMA) na liniach komunikacyjnych.
Pierwotnie oszacowano, że zaimplementowanie tej funkcjonalności nie
będzie konieczne, jednak przeprowadzone testy wykazały, że była to fałszywa hipoteza.